% book模板
% book模板.tex

\documentclass[12pt,UTF8]{ctexbook}

% 设置纸张信息。
\input{../../../../common/page}

% 设置字体，并解决显示难检字问题。
\xeCJKsetup{AutoFallBack=true}
% 注意：ctexbook类已默认设置SimSun为CJKrmdefault
% 以下设置用于确保字体回退和扩展字体可用
% 仅设置扩展字体以避免与默认设置冲突
\setCJKfamilyfont{hei}{SimHei}
\setCJKfamilyfont{kai}{KaiTi}

% 目录 chapter 级别加点（.）。
\usepackage{titletoc}
\titlecontents{chapter}[0pt]{\vspace{3mm}\bf\addvspace{2pt}\filright}{\contentspush{\thecontentslabel\hspace{0.8em}}}{}{\titlerule*[8pt]{.}\contentspage}

% 支持目录点击跳转
\usepackage[colorlinks,linkcolor=blue,citecolor=blue,urlcolor=blue]{hyperref}

% 设置 part 和 chapter 标题格式。
\ctexset{
	part/name= {第,卷},
	part/number={\chinese{part}},
	chapter/name={第,篇},
	chapter/number={\arabic{chapter}}
}

% 图片相关设置。
\usepackage{graphicx}
\graphicspath{{Images/}}

% 设置署名格式。
\newenvironment{shuming}{\hfill\zihao{4}}

% 注脚每页重新编号，避免编号过大。
\usepackage[perpage]{footmisc}

% 设置古文原文格式。
\newenvironment{yuanwen}{\bfseries\zihao{4}}

% 列表项向右偏移。
\usepackage{enumitem}

\title{\heiti\zihao{0} 电台}
\author{WangFei}
\date{\today}

\begin{document}

\maketitle
\tableofcontents

\frontmatter
\chapter{前言、序言}

电台是收音机收听的重要内容来源，本章节将详细介绍电台的相关知识，包括电台的分类、频率、呼号、节目内容以及收听技巧等内容。

\mainmatter

% 增加空行
~\\

\chapter{电台的基本概念}

\section{电台的定义}

\section{电台的分类}

\subsection{按服务对象分类}

\subsection{按广播方式分类}

\subsection{按发射功率分类}

\subsection{按调制方式分类}

\chapter{调幅电台}

\section{中波电台}

\subsection{中波频率范围}

\subsection{中波电台的特点}

\subsection{中波电台的传播特性}

\section{短波电台}

\subsection{短波频率范围}

\subsection{短波电台的特点}

\subsection{短波电台的传播特性}

\subsection{短波广播时间表}

\section{长波电台}

\subsection{长波频率范围}

\subsection{长波电台的特点}

\subsection{长波电台的传播特性}

\chapter{调频电台}

\section{调频频率范围}

\section{调频电台的特点}

\section{调频立体声}

\section{调频电台的传播特性}

\chapter{数字电台}

\section{数字广播技术}

\subsection{DAB数字广播}

\subsection{HD Radio}

\subsection{DMB数字多媒体广播}

\section{数字电台的特点}

\section{数字电台的发展趋势}

\chapter{国际广播电台}

\section{著名国际广播电台}

\subsection{BBC世界广播}

\subsection{VOA美国之音}

\subsection{NHK世界广播}

\subsection{RFA自由亚洲电台}

\subsection{AIR印度德里广播电台}

\subsection{HSK9-RT泰国广播电台}

\subsection{LNR老挝国际广播电台}

\subsection{PBC巴基斯坦广播电台}

\subsection{RAE阿根廷对外广播电台}

\subsection{RRI罗马尼亚国际广播电台}

\subsection{RRI-VOI印度尼西亚之声}

\subsection{RV梵蒂冈广播电台}

\subsection{TNVN-VOV越南之声}

\subsection{TRT-VOT土耳其之声}

\subsection{VOM蒙古之声}

\subsection{福音传教类电台}

\subsubsection{AWR-KSDA复临信徒世界广播网希望福音电台}

\subsubsection{FEBC远东广播公司}

\subsubsection{KNLS生命之光广播电台}

\subsubsection{KTWR环球广播电台}

\subsubsection{RVA天主教亚洲真理电台}

\subsection{其他国际广播电台}

\section{国际广播频率表}

\subsection{频率表使用说明}

使用说明：如无特别注明，本频率表有效期为2018年10月28日至2019年3月30日。本频率表中所标注之默认时间均为北京时间，默认频率单位皆为千赫兹，所列国际电台以英文字母缩写或全称排序。

免责声明：本频率表仅作电台频率数据客观记录及收听学习之用，请勿用于商业甚至非法用途。由于短波国际广播具有明显的政治倾向性，请收听者注意内容甄别！

\subsection{普通话国际广播频率表}

\subsubsection{AIR 印度德里广播电台}

\begin{itemize}[leftmargin=2cm]
    \item 频率：7250，7300，9655，11640，11770，13580，13655，15160
    \item 1000-1057：6180（古巴），6020（阿尔巴尼亚），7330，9570（阿尔巴尼亚），9825，11110，13655，13720，13750，15120，15160，17740
    \item 1500-1557：11710，11875，13750，17520，17650，17740
\end{itemize}

\subsubsection{HSK9-RT 泰国广播电台}

\begin{itemize}[leftmargin=2cm]
    \item 频率：7395，7430
\end{itemize}

\subsubsection{LNR 老挝国际广播电台}

\begin{itemize}[leftmargin=2cm]
    \item 2200-2230：567，6130
\end{itemize}

\subsubsection{PBC 巴基斯坦广播电台（A18，发射机长期损坏）}

\begin{itemize}[leftmargin=2cm]
    \item 2000-2100：11590，15730
\end{itemize}

\subsubsection{RAE 阿根廷对外广播电台（A18）}

\begin{itemize}[leftmargin=2cm]
    \item 1700-1800：5950（周二至周六）
\end{itemize}

\subsubsection{RRI 罗马尼亚国际广播电台}

\begin{itemize}[leftmargin=2cm]
    \item 2130-2200：9610，11825
    \item 1300-1330：11820，13730（DRM）
\end{itemize}

\subsubsection{RRI-VOI 印度尼西亚之声}

\begin{itemize}[leftmargin=2cm]
    \item 1900-2000：3325
    \item 2300-0000：3325
\end{itemize}

\subsubsection{RV 梵蒂冈广播电台}

\begin{itemize}[leftmargin=2cm]
    \item 0600-0630：6185，7410，9580
    \item 2030-2115（周六）：6115，7485，9560
\end{itemize}

\subsubsection{TNVN-VOV 越南之声}

\begin{itemize}[leftmargin=2cm]
    \item 0600-0630：7220
    \item 0630-0700：9840，12020
    \item 1900-1930：7220
    \item 2000-2030：7220
    \item 2100-2130：7220
\end{itemize}

\subsubsection{TRT-VOT 土耳其之声}

\begin{itemize}[leftmargin=2cm]
    \item 2000-2100：13630
\end{itemize}

\subsubsection{VOM 蒙古之声}

\begin{itemize}[leftmargin=2cm]
    \item 1800-1830：12085
    \item 2230-2300：12015
\end{itemize}

\subsubsection{福音传教类电台}

\paragraph{AWR-KSDA 复临信徒世界广播网希望福音电台（A18）}

\begin{itemize}[leftmargin=2cm]
    \item 频率：15715，15285（周一至周五）
\end{itemize}

\paragraph{FEBC 远东广播公司}

\subparagraph{良友电台·第一台}

\begin{itemize}[leftmargin=2cm]
    \item 0700-0900：12070
    \item 1830-2200：9400
    \item 2200-0000：9345
\end{itemize}

\subparagraph{良友电台·第三台}

\begin{itemize}[leftmargin=2cm]
    \item 0630-0830：9405
\end{itemize}

\subparagraph{良友电台·iRadio 爱广播频道}

\begin{itemize}[leftmargin=2cm]
    \item 1800-0000：9275
\end{itemize}

\subparagraph{HLAZ 益友电台·第一台}

\begin{itemize}[leftmargin=2cm]
    \item 1900-2030：1566
    \item 2145-0030：1566
\end{itemize}

\paragraph{KNLS 生命之光广播电台}

\subparagraph{阿拉斯加发射}

\begin{itemize}[leftmargin=2cm]
    \item 1600-2000：7355
    \item 2100-2200：7355，7370
    \item 2200-0200：7355
\end{itemize}

\subparagraph{马达加斯加发射}

\begin{itemize}[leftmargin=2cm]
    \item 1200-1300：17530（华南方向）
    \item 0500-0600：11610（欧洲方向）
\end{itemize}

\paragraph{KTWR 环球广播电台}

\begin{itemize}[leftmargin=2cm]
    \item 1830-1900：12120
    \item 1900-2000：9910（周一至周五）
    \item 1930-2000：9910（周日）
    \item 1945-2000：9975
    \item 2000-2015：9910（周日至周五）
    \item 2115-2130：9975（周六）
    \item 2115-2145：9975（周一至周五）
    \item 2145-2150：9975（周一至周五）
\end{itemize}

\paragraph{RVA 天主教亚洲真理电台（退出短波）}

\section{国际广播的作用}

\section{国际广播的收听特点}

\chapter{电台呼号与频率}

\section{电台呼号的组成}

\section{呼号的分配规则}

\section{频率分配与管理}

\section{电台的识别方法}

\chapter{电台节目内容}

\section{新闻类节目}

\section{音乐类节目}

\section{访谈类节目}

\section{教育类节目}

\section{娱乐类节目}

\section{专题类节目}

\chapter{电台收听技巧}

\section{频率查找方法}

\subsection{使用频率表}

\subsection{使用网络资源}

\subsection{使用移动应用}

\section{最佳收听时间}

\subsection{中波最佳收听时间}

\subsection{短波最佳收听时间}

\subsection{调频最佳收听时间}

\section{天线选择与架设}

\section{干扰排除方法}

\section{远程接收技巧}

\chapter{电台日志与QSL卡片}

\section{电台日志的记录}

\section{QSL卡片的意义}

\section{QSL卡片的申请方法}

\section{QSL卡片的收集与展示}

\chapter{业余无线电台}

\section{业余电台概述}

\section{业余电台的频段}

\section{业余电台的操作模式}

\section{业余电台的呼号分配}

\section{业余电台的执照要求}

\backmatter

\end{document}
